% =============================================================================
% CAPITULO 7 - DISCUSION
% =============================================================================
\chapter{Discusión}
\label{ch:discusion}

Este capítulo discute los resultados del capítulo~\ref{ch:resultados}
desde tres ángulos sucesivos: qué demuestran los datos del Mundial
Qatar 2022 sobre la pregunta de investigación; qué no demuestran y
por qué la lectura debe ser cautelosa en varios frentes; y dónde el
\CCV funciona como herramienta de medida y dónde su poder es
insuficiente. El capítulo cierra con las implicaciones operativas
para \textit{scouting}, mercado de fichajes, preparación de partidos
de fase eliminatoria y gestión de cambios desde el banquillo.


\section{Qué demuestran los resultados}
\label{sec:disc-demuestran}

El primer hallazgo es que la señal \textit{clutch} a nivel individual
es identificable causalmente sobre un corpus relativamente pequeño
---ciento setenta y dos \textit{shocks} de gol en un torneo
eliminatorio--- cuando la identificación se hace condicional al
contexto adecuado. La dimensión sobre la que la identificación es
estable es la respuesta a la proximidad de eliminación: tres
jugadores muestran efecto positivo significativo a probabilidad
posterior $\geq 0{,}85$ ---Mbappé, Rashford y Saka--- y uno muestra
efecto negativo significativo ---Muntari---, todos consistentes con
la narrativa cualitativa de sus partidos en el torneo. La magnitud
del efecto de Mbappé ($+0{,}110$ desviaciones estándar del canal con
probabilidad posterior $0{,}97$) sobre la final Argentina-Francia se
alinea con la observación pública de su irrupción concentrada en los
últimos veinte minutos de regulación y la prórroga. El marco recupera
ese caso como una propiedad medida del jugador, no como una
intuición \textit{post hoc} del observador.

El segundo hallazgo invierte la lectura habitual del ATE. La
distribución del CATE individual reportada en
§\ref{sec:res-cate} ---$\sigma = 0{,}086$ en ofensivo tras marcar,
rango $[-0{,}49,\, +0{,}37]$ desviaciones estándar--- es
sustantivamente distinta de cero, mientras que las ocho celdas
canal-por-contexto del ATE poblacional son indistinguibles de cero.
La consecuencia operativa para el scouting es directa: un ATE nulo
no implica ausencia de efecto, implica que los efectos individuales
de signo opuesto se cancelan al agregarlos. La firma corresponde a
la concepción clásica del \textit{clutch} en deporte
\parencite{Otten2009Clutch}: el mismo \textit{shock} produce
respuestas opuestas en jugadores distintos, y el promedio cancela
esa variabilidad.

El tercer hallazgo es la coherencia parcial con la valoración humana
experta. De los cinco canales-índice contrastados contra el
\textit{grade} PFF agregado, dos resultan significativos
(\textit{off-ball} $\rho = +0{,}211$ y Remontador $\rho = +0{,}177$,
detalle en §\ref{sec:res-validacion-pff}) y tres no. La lectura es
que el marco recoge una dimensión reconocible por evaluadores
humanos en parte de los canales, pero CATE y \textit{grades} no son
medidas intercambiables: el \textit{grade} agrega el partido en un
único valor por jugador, mientras el CATE residualiza por minuto y
por \textit{shock}.

El cuarto hallazgo es metodológico. La identificación causal a nivel
individual se sostiene sobre la composición de las cinco capas
formalizadas en el capítulo~\ref{ch:metodologia}: el contraste
temporal pre/post sobre el mismo jugador, los efectos fijos
jugador-\textit{shock} en el DiD, la residualización
\textit{leave-one-out} contra el bloque del propio equipo, el
cuasi-experimento \textit{near-miss} y el AIPW con doble robustez.
El cuasi-experimento de \textit{near-miss} aporta una segunda vía
identificativa sobre el corpus de 70 disparos pre-registrados; con
ese tamaño, el brazo opera a nivel agregado y no permite
verificación cruzada por jugador, restricción declarada en
§\ref{sec:disc-no-demuestran}. El test $F$ de pre-tendencias
agregadas \parencite{Roth2022Pretest} no rechaza la hipótesis de
paralelismo en ninguno de los ocho contrastes, y las cotas
\textit{HonestDiD} \parencite{RambachanRoth2023} para los tres
niveles de relajación contienen el cero en todos los casos. Los
coeficientes nulos del ATE no parecen, por tanto, un artefacto de
violación de identificación, sino una propiedad del corpus.


\section{Qué no demuestran los resultados}
\label{sec:disc-no-demuestran}

El primer punto de cautela es la interpretación del ATE poblacional.
Que los ocho coeficientes canal-por-contexto sean indistinguibles de
cero al cinco por ciento no es prueba de que el \textit{shock} no
mueva al jugador promedio. Es prueba de que, con el tamaño de muestra
disponible ---en torno a cincuenta \textit{shocks} efectivos por
celda tras filtros de ventana y disponibilidad---, el efecto promedio
no es detectable por encima del nivel de ruido. El análisis de
potencia ex-ante sitúa el \textit{Minimum Detectable Effect} (MDE) al
ochenta por ciento de potencia entre $0{,}0003$ desviaciones estándar
para el canal \textit{off-ball} y $0{,}036$ para el canal físico. Si
el verdadero efecto poblacional cae por debajo de esos umbrales, el
marco no lo recupera con $N = 172$ \textit{shocks}; lo que se puede
decir es que el efecto promedio no es lo bastante grande para
discernirse sobre este corpus, no que sea cero. Por eso el capítulo
previo reporta el ATE con su IC y declara la indistinguibilidad de
cero en lugar de afirmar la nulidad.

El segundo punto es la generalización fuera del torneo. La aplicación
del \CCV se restringe al Mundial Qatar 2022 porque es el único corpus
público que combina los tres ingredientes estructurales que el marco
exige simultáneamente: \textit{eventing} sincronizado con tracking
continuo \textit{full-pitch} de los 22 jugadores a 25\,Hz y cobertura
completa de fase eliminatoria. PFF añade además
\textit{grades} humanos por evento, que el marco aprovecha como
\textit{prior} informativo y validación externa pero no exige como
condición de aplicabilidad. Las propiedades individuales
identificadas para los cuatro jugadores con efecto significativo
\textit{pressure-clutch} corresponden a su comportamiento en este
torneo concreto; no son una predicción de su comportamiento en otros
contextos competitivos ni una validación de su perfil ``\textit{clutch}''
como rasgo estable. Replicar el marco sobre la Eurocopa 2024, sobre
los partidos de eliminatoria de la \textit{Champions League} o sobre
cualquier otro torneo abierto futuro permitiría empezar a hablar de
la estabilidad temporal de la propiedad; el corpus actual solo
permite observar su existencia en un único torneo.

El tercer punto es la limitación del brazo del cuasi-experimento de
\textit{near-miss}. El corpus de 70 \textit{near-misses} sobre las
cinco tipologías es reducido y deja al AIPW sin potencia para
detectar efectos individuales. La triangulación entre DiD y
\textit{near-miss} se reporta a nivel agregado y coincide
cualitativamente en signo, pero a nivel individual el contraste entre
los dos brazos no es testable con este corpus. Por eso el trabajo
toma el DiD \textit{within-player} como brazo principal y reporta el
AIPW como verificación complementaria; la triangulación rigurosa por
jugador queda como línea futura cuando los corpus disponibles
permitan ejecutarla.

El cuarto punto es que el \CCV mide una forma específica de
\textit{clutch}, no \textit{clutch} en sentido amplio. Los tres
contextos del marco ---post-gol a favor, post-gol en contra y
proximidad de eliminación alta--- identifican momentos donde la carga
emocional y la presión competitiva cruzan un umbral observable a
nivel de evento o de contexto. Hay situaciones de presión competitiva
genuina que quedan fuera del marco por no encajar en ese umbral: un
mediocampista que protege un empate en el minuto noventa de un
partido de fase de grupos con ambos equipos ya clasificados está bajo
presión, pero sin un \textit{shock} emocional discreto ni proximidad
de eliminación inminente que el marco pueda etiquetar como
tratamiento. La identificación causal a nivel individual exige momentos
con umbral identificable; el contraste de un jugador en momentos
\textit{clutch} difusos sin ese umbral queda fuera del alcance del
trabajo. La frontera no es metodológica sino de muestra: con un
corpus mayor y un catálogo más amplio de tipologías de
\textit{trigger} ---por ejemplo, segmentos de partido condicionados
al estado del marcador combinado con la trascendencia clasificatoria,
o ventanas en torno a sustituciones tácticas---, el alcance operativo
del marco podría extenderse a momentos hoy fuera del umbral
identificable.

El quinto punto matiza el alcance de la correlación con los
\textit{grades} humanos PFF. Los dos coeficientes significativos
(reportados arriba en §\ref{sec:disc-demuestran}) se calculan sobre
166 jugadores con muestra fiable, un tercio del corpus total. La
interpretación adecuada no es que los CATE midan ``lo mismo'' que el
ojo humano experto: el CATE estima el efecto del \textit{shock}
residualizado contra el bloque, y el \textit{grade} agrega el
rendimiento completo del partido. El alineamiento es direccional,
no de magnitud.


\section{Dónde funciona y dónde no funciona el \CCV}
\label{sec:disc-donde-funciona}

El \CCV resulta útil como herramienta de medida en tres frentes
sobre el corpus disponible. El primero es la dimensión
\textit{pressure-clutch}: la modulación continua del CATE por la
proximidad de eliminación produce efectos individuales estables y
medibles con la muestra disponible. Cuatro jugadores cruzan el
umbral de significatividad a probabilidad posterior $\geq 0{,}85$,
todos con narrativa cualitativa coherente con sus partidos en el
torneo. El segundo es la celda canal-específica
ataque-tras-marcar: veintidós jugadores presentan efecto
significativo, repartidos entre cinco con respuesta positiva y
diecisiete con respuesta negativa. La presencia de Mbappé en ambos
grupos ---positivo en \textit{pressure-clutch} y negativo en
ataque-tras-marcar--- es informativa: el marco distingue, sobre el
mismo jugador, propiedades contextualmente diferenciadas que un único
agregado escondería. El tercer frente es la lectura
\textit{scout-facing}: la tabla canónica del \CCV, con un jugador por
fila y los tres índices condensados, es operativa para uso de
\textit{scouting} sobre el corpus.

El marco no funciona, sobre este corpus, en los agregados Remontador
y Cerrojo. La razón es estructural y no metodológica. Cada agregado
es la suma ponderada de cuatro celdas canal-por-contexto cuyas
posteriores no están correlacionadas en signo entre canales, de
modo que el efecto agregado se suaviza al promediarlas. Con $N = 172$
\textit{shocks} repartidos entre 120 de fase de grupos y 52 de
eliminatoria, el agregado por dirección queda infra-identificado
incluso para los jugadores con la respuesta individual más fuerte.
La identificación condicional al contexto de eliminación
inminente, en cambio, opera sobre un eje único de modulación y
permite que el efecto individual emerja sobre el ruido.

El marco tampoco funciona, sobre este corpus, en el canal físico a la
ventana de referencia $\pm 10$\,minutos. El análisis de sensibilidad
reportado en §\ref{sec:res-identificacion} sí detecta efectos
significativos al ampliar a $\pm 15$\,minutos, lo que sugiere una
respuesta física diferida fuera de la ventana inmediata posterior al
\textit{shock}. La ventana $\pm 10$ se mantiene como referencia por
su correspondencia con la duración del \textit{shock} emocional en la
psicología del deporte; el canal físico a horizonte ampliado queda
como línea futura.

\section{Implicaciones operativas}
\label{sec:disc-implicaciones}

\paragraph{Scouting y mercado de fichajes.}
El \CCV produce una tabla \textit{scout-facing} con un jugador por
fila y tres índices ---Remontador, Cerrojo y \textit{Pressure
Response}--- acompañados de la media posterior, el intervalo de
credibilidad y la probabilidad posterior por jugador, además de los
\textit{rankings} dentro de su posición granular y dentro de su
\textit{bucket} posicional. Esa estructura permite separar, para un
fichaje concreto, la prima clutch atribuible al jugador de la
atribuible al bloque del equipo en el que ha rendido históricamente.
Un director deportivo que valora a un jugador con reputación
\textit{clutch} adquirida en un equipo grande puede inspeccionar el
delta relativo al bloque LOO sobre el corpus disponible y comprobar
si la propiedad sobrevive al control por el patrón colectivo de su
equipo en cada \textit{shock}. La interpretación recomendada es
cautelosa: el marco no produce una decisión, produce una métrica
informada por evidencia causal que se combina con el conocimiento
táctico y económico del propio departamento.

\paragraph{Preparación de partidos de fase eliminatoria.}
La identificación de \textit{pressure-clutch} a nivel individual
sobre el contexto de eliminación inminente tiene implicaciones
directas para la preparación de partidos KO. Los jugadores con
efecto positivo significativo bajo proximidad de eliminación alta
son candidatos naturales a recibir más minutos en los partidos de
mayor presión, mientras que los jugadores con efecto negativo
significativo son candidatos a ser rotados antes de los tramos
críticos del calendario. La decisión sigue siendo del cuerpo
técnico; el marco proporciona una métrica con identificación causal
sobre la respuesta histórica del jugador a este contexto concreto,
no una predicción individual sobre el siguiente partido.

\paragraph{Gestión de cambios y rotación.}
La celda ataque-tras-marcar muestra que diecisiete jugadores caen
significativamente en su producción ofensiva inmediata después de
anotar, incluido Mbappé. Algunos perfiles ofensivos aflojan en los
minutos posteriores al gol, justo cuando el equipo es más vulnerable
a un contragolpe; para esos perfiles concretos puede tener sentido
rotar o ajustar tácticamente la ventana post-marcador. Los cinco
jugadores con respuesta positiva en la misma celda ---Choupo-Moting,
Vlahović, Sterling, Kane y Cristiano Ronaldo--- presentan el patrón
opuesto: mantener su peso ofensivo después del gol empuja la
búsqueda del segundo tanto.

\paragraph{Lectura general y cuándo no usar el marco.}
El \CCV es una herramienta de medida bajo supuestos identificables,
no un oráculo sobre la propiedad \textit{clutch} de un jugador. Produce
estimaciones con intervalo de credibilidad y probabilidad posterior;
la decisión final siempre suma información que el marco no observa
---táctica, psicológica, médica, contextual del partido siguiente.
Hay tres usos en los que el marco no debe sustituir al juicio experto:
la valoración de jugadores con muestra \textit{low\_sample} (marcados
explícitamente en la tabla), la proyección de la propiedad a contextos
competitivos no observados y la atribución \textit{post hoc} del
resultado de un partido a un \textit{shock} concreto. En esos usos
el marco aporta información, pero no concluye; el reporte sigue siendo
la media posterior con su intervalo, no un valor puntual.
